\documentclass[11pt]{article}

\usepackage[margin=1in]{geometry}
\usepackage{amsmath,amssymb,mathtools}
\usepackage{booktabs}
\usepackage{array}
\usepackage{hyperref}

\newcommand{\SU}{\mathrm{SU}}
\newcommand{\Spin}{\mathrm{Spin}}
\newcommand{\Tr}{\mathrm{Tr}}
\newcommand{\Res}{\mathrm{Res}}
\newcommand{\diag}{\mathrm{diag}}
\newcommand{\half}{\frac{1}{2}}
\newcommand{\bml}{B-L}
\newcommand{\uc}{u^c}
\newcommand{\dc}{d^c}
\newcommand{\nuc}{\nu^c}
\newcommand{\ec}{e^c}

\title{Route C Bootstrap Derivation Ledger through P8}
\author{BoYu Chen}
\date{\today}

\begin{document}
\maketitle

\begin{abstract}
This note records the derivations generated in Route C stages P1--P8.
It is a derivation ledger, not a claim that on-shell consistency has already
derived a complete GUT.  The current Route C program uses the one-family
Standard Model spectrum with a right-handed neutrino as baseline, builds
charge-filtered symbolic four-point pole components, checks conditional
residue factorization/positivity, marks Ward-identity completion requirements,
compares anomaly/chiral consistency across $\SU(5)$, Pati--Salam,
$\Spin(10)$, $E_6$, and superstring-derived completions, and classifies
controlled high-energy-growth and low-energy proton-bound readiness
requirements, and closes the first pass with a comparative theory scorecard.
\end{abstract}

\section{Boundary of the Ledger}

Route C tests whether
\begin{equation}
\begin{split}
  &\text{Lorentz}
  + \text{unitarity}
  + \text{factorization}
  + \text{Ward identities} \\
  &\qquad
  + \text{anomaly cancellation}
  + \text{controlled high-energy behavior}
\end{split}
\end{equation}
can filter conditional GUT branches.  It must not be read as
\begin{equation}
  \text{four axioms} \Longrightarrow \text{the full GUT}.
\end{equation}
The honest statement is instead
\begin{equation}
  \text{on-shell consistency} \Longrightarrow
  \text{filters on allowed action-level GUT branches}.
\end{equation}

The baseline external spectrum is one left-handed family plus a
right-handed-neutrino conjugate field,
\begin{equation}
  \mathcal F_{\rm SM+\nu^c}
  =
  Q \oplus L \oplus \uc \oplus \dc \oplus \nuc \oplus \ec .
\end{equation}
All fields are written as left-handed Weyl fields.

\section{P1: External-State Charges and Anomalies}

\subsection{Charge Table}

The P1 charge ledger fixes the Standard Model face conventions shown in
Table~\ref{tab:p1-charges}.

\begin{table}[h]
\centering
\begin{tabular}{c c c c c c}
\toprule
state & $\SU(3)_C$ & $\SU(2)_L$ & $Y$ & $\bml$ & multiplicity \\
\midrule
$Q$    & $3$        & $2$ & $1/6$   & $1/3$  & $6$ \\
$L$    & $1$        & $2$ & $-1/2$  & $-1$   & $2$ \\
$\uc$  & $\bar 3$   & $1$ & $-2/3$  & $-1/3$ & $3$ \\
$\dc$  & $\bar 3$   & $1$ & $1/3$   & $-1/3$ & $3$ \\
$\nuc$ & $1$        & $1$ & $0$     & $1$    & $1$ \\
$\ec$  & $1$        & $1$ & $1$     & $1$    & $1$ \\
\bottomrule
\end{tabular}
\caption{P1 baseline spectrum.  The fields $\uc,\dc,\nuc,\ec$ are
left-handed conjugate Weyl fields, so the listed charges are the charges
of those conjugate fields.}
\label{tab:p1-charges}
\end{table}

\subsection{Hypercharge Normalization}

For one full family plus $\nuc$, the hypercharge trace is
\begin{align}
  \Tr Y^2
  &=
  6\left(\frac{1}{6}\right)^2
  +2\left(-\frac{1}{2}\right)^2
  +3\left(-\frac{2}{3}\right)^2
  +3\left(\frac{1}{3}\right)^2
  +1\cdot 0^2
  +1\cdot 1^2                                              \notag\\
  &=\frac{10}{3}.
\end{align}
With $\Tr_{\mathbf 2}T_{3L}^2=1/2$, the weak trace is
\begin{equation}
  \Tr T_{3L}^2
  =
  3\cdot \frac{1}{2}
  +1\cdot \frac{1}{2}
  =2.
\end{equation}
Therefore
\begin{equation}
  \frac{\Tr Y^2}{\Tr T_{3L}^2}=\frac{5}{3},
\end{equation}
which is the standard GUT hypercharge normalization.

\subsection{Perturbative and Global Anomaly Checks}

With the conventions of Table~\ref{tab:p1-charges}, the exact rational
ledger verifies
\begin{equation}
  \mathcal A_{\SU(3)^3}
  =
  2A(3)+A(\bar 3)+A(\bar 3)
  =
  2-1-1=0 .
\end{equation}
The mixed color-hypercharge anomaly is
\begin{equation}
  \mathcal A_{\SU(3)^2U(1)_Y}
  =
  2T(3)\frac{1}{6}
  +T(\bar 3)\left(-\frac{2}{3}\right)
  +T(\bar 3)\frac{1}{3}
  =
  \frac{1}{6}-\frac{1}{3}+\frac{1}{6}=0,
\end{equation}
where $T(3)=T(\bar 3)=1/2$.  The weak-hypercharge anomaly is
\begin{equation}
  \mathcal A_{\SU(2)^2U(1)_Y}
  =
  3T(2)\frac{1}{6}
  +T(2)\left(-\frac{1}{2}\right)
  =
  \frac{1}{4}-\frac{1}{4}=0.
\end{equation}
The gravitational and cubic hypercharge anomalies are
\begin{align}
  \mathcal A_{\mathrm{grav}^2U(1)_Y}
  &=
  6\frac{1}{6}
  +2\left(-\frac{1}{2}\right)
  +3\left(-\frac{2}{3}\right)
  +3\frac{1}{3}
  +0+1
  =0,                                                        \\
  \mathcal A_{U(1)_Y^3}
  &=
  6\left(\frac{1}{6}\right)^3
  +2\left(-\frac{1}{2}\right)^3
  +3\left(-\frac{2}{3}\right)^3
  +3\left(\frac{1}{3}\right)^3
  +0^3+1^3
  =0.
\end{align}
The same exact ledger verifies the analogous $\bml$ and mixed
$Y$--$\bml$ anomalies,
\begin{equation}
  \SU(3)^2(\bml),\quad
  \SU(2)^2(\bml),\quad
  \mathrm{grav}^2(\bml),\quad
  (\bml)^3,\quad
  Y^2(\bml),\quad
  Y(\bml)^2,
\end{equation}
all vanish.  The Witten $\SU(2)$ check passes because one family contains
three $Q$ doublets and one $L$ doublet, hence four left-handed doublets.

\section{P2: Symbolic Four-Point Pole Ansatz}

The first pole ansatz is an all-incoming chiral-bilinear seed,
\begin{equation}
  \mathcal A_4(12|34)
  =
  \sum_X\frac{R_s(X)}{s-M_X^2}+C_6(12|34),
  \qquad
  R_s(X)=g_{12X}g_{34\bar X}.
\end{equation}
P2 does not declare the mediator spectrum.  It only computes which
symbolic pole components are not forbidden by the P1 charges and by
non-Abelian conjugacy of pair products.

The generated P2 ledger gives the numerical summary
\begin{equation}
\begin{array}{c|c}
\text{quantity} & \text{value} \\
\hline
\text{pair types} & 21\\
\text{pair-product components} & 31\\
\text{tested pair-component pairings} & 496\\
\text{allowed symbolic pole components} & 18\\
\text{forbidden by Abelian charge} & 455\\
\text{forbidden after Abelian pass by non-Abelian conjugacy} & 23
\end{array}
\end{equation}
The proton-operator seeds retained at this purely symbolic level are
\begin{equation}
  QQQL:\ 2 \ \text{allowed pole components},\qquad
  \uc\uc\dc\ec:\ 2 \ \text{allowed pole components}.
\end{equation}
The entries
\begin{equation}
  QQ\uc\ec,\qquad QL\uc\dc
\end{equation}
have zero allowed pole components in this first all-incoming symbolic
seed.  These zeros are not no-go theorems for full models; they are
bookkeeping results for the declared P2 ansatz and charge conventions.

\section{P3: Residue Factorization and Positivity}

For a declared mediator sector $X$, let $\alpha$ label all two-particle
channels that can couple to $X$.  Near a positive-norm simple pole,
tree-level factorization requires
\begin{equation}
  \Res_X \mathcal A_{\alpha\beta}
  =
  g_{\alpha X}g_{\beta X}^\ast .
\end{equation}
Thus the residue matrix is a Gram matrix,
\begin{equation}
  R_X = g_X g_X^\dagger,
\end{equation}
and is positive semidefinite in a positive-metric Hilbert space:
\begin{equation}
  v^\dagger R_X v = |g_X^\dagger v|^2\ge 0.
\end{equation}
A wrong-sign mediator metric reverses the sign and fails the unit-coupling
stress test.

The P3 machine ledger checks the P2-allowed poles and reports
\begin{equation}
\begin{array}{c|c}
\text{quantity} & \text{value} \\
\hline
\text{mediator sectors} & 26\\
\text{total vertex channels} & 31\\
\text{P2 allowed poles checked} & 18\\
\text{factorization checks passed} & 18\\
\text{PSD blocks under positive metric} & 26\\
\text{largest channel-block size} & 3
\end{array}
\end{equation}
For a unit-coupling block with $n$ channels, the stress spectrum is
\begin{equation}
  \mathrm{spec}(R_X)=\{n,0,\ldots,0\},
  \qquad
  \mathrm{spec}(-R_X)=\{-n,0,\ldots,0\},
\end{equation}
which is the numerical pattern seen in the P3 output.

\section{P4: Ward-Identity Bookkeeping}

For an unbroken massless gauge boson, the target Ward check is
\begin{equation}
  \epsilon_\mu\rightarrow p_\mu,\qquad p_\mu\mathcal M^\mu=0.
\end{equation}
For a broken massive vector, the target relation is the generalized
Ward/Slavnov--Taylor form
\begin{equation}
  p_\mu\mathcal M^\mu(V_X)=M_X\mathcal M(\phi_X),
\end{equation}
so a Goldstone/Higgs/source completion is required before a true Ward
identity proof can be claimed.

The P4 ledger separates unbroken SM currents from broken/off-face
currents:
\begin{equation}
\begin{array}{c|c}
\text{quantity} & \text{value} \\
\hline
\text{unbroken SM current sectors} & 3\\
\text{current transitions checked} & 36\\
\text{massless diagonal transitions ready} & 6\\
\text{broken/off-face transitions needing completion} & 30\\
\text{mediator sectors inspected} & 26\\
\text{mediator sectors needing completion if vector} & 26
\end{array}
\end{equation}
The bookkeeping-ready unbroken sectors are
\begin{equation}
  \SU(3)_C,\qquad \SU(2)_L,\qquad U(1)_Y.
\end{equation}
All off-face $\Spin(10)$, Pati--Salam leptoquark, and $\SU(2)_R$
charged-current transitions are flagged as needing explicit generator
matrices, massive-vector data, and Goldstone/Higgs/source completion.

\section{P5: Anomaly and Chiral-Consistency Candidate Ledger}

P5 compares candidate branches against the P1 baseline.  The structural
filter used for the minimal Route C field-theory comparison is
\begin{equation}
\begin{split}
  &\text{simple group}
  +\text{one irreducible family object}
  +\text{contains }\nuc  \\
  &\qquad
  +\text{no chiral SM exotics}
  +\text{minimal standard realization}.
\end{split}
\end{equation}

\subsection{SU(5)}

One family plus $\nuc$ is realized as
\begin{equation}
  10\oplus \bar 5\oplus 1.
\end{equation}
The cubic anomaly cancels:
\begin{equation}
  A(10)+A(\bar 5)+A(1)=1-1+0=0.
\end{equation}
Also $\pi_4(\SU(5))=0$ in the standard four-dimensional global-anomaly
check.  However, the family is not one irreducible object; therefore
$\SU(5)$ is anomaly-consistent but fails the single-object filter.

\subsection{Pati--Salam}

The standard Pati--Salam family is
\begin{equation}
  (4,2,1)\oplus(\bar 4,1,2)
  \quad\text{under}\quad
  \SU(4)_C\times\SU(2)_L\times\SU(2)_R .
\end{equation}
The $\SU(4)_C^3$ anomaly cancels because
\begin{equation}
  2A(4)+2A(\bar 4)=2-2=0.
\end{equation}
The Witten checks pass because $\dim(4)=4$ gives four $\SU(2)_L$
doublets and four $\SU(2)_R$ doublets.  Pati--Salam is physically natural
and anomaly-consistent, but it is a product group and not a single
irreducible family object under a simple group.

\subsection{Spin(10)}

The standard $\Spin(10)$ family is the half-spinor
\begin{equation}
  16.
\end{equation}
It contains the full
\begin{equation}
  Q\oplus L\oplus\uc\oplus\dc\oplus\nuc\oplus\ec
\end{equation}
after restriction to the Standard Model face.  The $D_5$ algebra has no
independent cubic gauge-anomaly invariant in this standard check, and
$\pi_4(\Spin(10))=0$.  Therefore $\Spin(10):16$ is anomaly-safe and is the
standard minimal single-object survivor under the stated Route C
pre-filter.

\subsection{E6}

The $E_6$ fundamental representation restricts as
\begin{equation}
  27\longrightarrow 16\oplus 10\oplus 1
  \qquad \text{under } \Spin(10).
\end{equation}
It contains a $\Spin(10)$ family, but also contains extra $10\oplus 1$
states.  Since $E_6$ has no independent cubic gauge-anomaly invariant for
the standard $27$, the anomaly status is acceptable, but the branch
requires an explicit lifting/vectorlike audit for the extra states.  It is
therefore not the minimal Route C survivor.

\subsection{Superstring-Derived Completions}

Superstring-derived completions are not a single four-dimensional
field-theory group.  They form a UV-completion class.  A string-derived
branch must supply compactification-specific data:
\begin{equation}
  \text{Green--Schwarz/tadpole/modular consistency},
  \qquad
  \text{massless spectrum},
  \qquad
  \text{chiral index},
\end{equation}
together with exotics lifting and matching to the Route-A
representation/index skeleton.  Such branches may provide stronger
high-energy softness or tower structure, but P5 does not treat them as a
single fixed four-dimensional candidate.

\subsection{P5 Numerical Summary}

The generated P5 ledger reports
\begin{equation}
\begin{array}{c|c}
\text{quantity} & \text{value}\\
\hline
\text{candidates inspected} & 5\\
\text{field-theory candidates} & 4\\
\text{UV-completion classes} & 1\\
\text{P1 baseline anomalies pass} & \text{true}\\
\text{minimal single-object survivor} & \Spin(10)
\end{array}
\end{equation}
This result should be read narrowly: passing anomaly checks alone does not
select a GUT.  The selection of $\Spin(10)$ occurs only after imposing the
Route C minimal single-object family filter.

\section{P6: Controlled High-Energy-Growth Audit}

P6 consumes the P3 residue blocks and the P4 Ward-completion ledger.  Its
purpose is not to compute a full four-point amplitude.  Its purpose is to
classify which symbolic branches would have uncontrolled high-energy growth
if treated as isolated finite broken-vector exchange.

For a finite massive vector, the dangerous contribution is the longitudinal
propagator term.  A conservative schematic form is
\begin{equation}
  \mathcal A_L
  \sim
  \frac{(J\cdot p)(J'\cdot p)}{M_X^2(s-M_X^2)}.
\end{equation}
If the broken current is not completed by the corresponding Goldstone/Higgs
sector, this can scale as
\begin{equation}
  \mathcal A_L \sim \frac{s}{M_X^2}
  \quad\text{or worse}
\end{equation}
in the high-energy regime.  Therefore the required broken-vector target is
again the generalized Ward/Slavnov--Taylor identity
\begin{equation}
  p_\mu\mathcal M^\mu(V_X)=M_X\mathcal M(\phi_X).
\end{equation}

The generated P6 ledger reports
\begin{equation}
\begin{array}{c|c}
\text{quantity} & \text{value}\\
\hline
\text{mediator sectors inspected} & 26\\
\text{finite-vector branches needing completion} & 26\\
\text{unbroken-vector branches bookkeeping-ready} & 0\\
\text{scalar/auxiliary branches without longitudinal-vector risk} & 26\\
\text{contact-only components needing UV completion} & 18\\
\text{broken transitions needing completion} & 30\\
\text{tower-like UV forced now} & \text{false}
\end{array}
\end{equation}
Thus every current P3 mediator sector, if interpreted as a finite broken
vector, needs explicit finite-completion data:
\begin{equation}
  \{\text{broken generator},\ M_X,\ \phi_X,\ \text{Higgs/source sector},
  \ \text{symmetry-fixed contact/numerator terms}\}.
\end{equation}

If a symbolic mediator is instead interpreted as scalar or auxiliary
exchange, the longitudinal-vector Ward problem is absent.  This is only a
bookkeeping statement; spin-statistics, action existence, low-energy
matching, and perturbative unitarity remain separate checks.

A contact-only four-fermion EFT normalization generically scales like
\begin{equation}
  \mathcal A_{\rm contact}\sim \frac{s}{\Lambda^2},
\end{equation}
so contact terms may encode low-energy matching but cannot by themselves
serve as high-energy softness proofs.

Finally, P6 does not force a string or tower completion.  It states the
conditional comparison rule:
\begin{equation}
\begin{split}
  \text{finite Goldstone/Higgs/source completion fails}
  \quad\Longrightarrow\quad\\
  \text{tower-like UV completion becomes a serious comparison class}.
\end{split}
\end{equation}
Superstring-derived completions therefore remain comparison branches, not
P6 conclusions.

\section{P7: Low-Energy Matching and Proton-Bound Readiness}

P7 converts the P2/P3 symbolic pole data into conditional low-energy
dimension-six matching rows.  It respects the P6 gate: a vector interpretation
of a broken mediator cannot be promoted to a physical Wilson coefficient until
the finite Goldstone/Higgs/source completion has been supplied.

For a symbolic pole with mediator $X$, P7 records the matching template
\begin{equation}
  C_6[\mathcal O;X]
  =
  \frac{g_{\mathrm{left},X}g_{\mathrm{right},X}^\ast}{M_X^2}.
\end{equation}
This is only a gauge-basis symbolic coefficient.  It is not yet a physical
proton-decay Wilson coefficient.

Using the left-handed Weyl charge convention, the global bookkeeping assigns
\begin{equation}
  B(Q)=\frac13,\quad B(\uc)=B(\dc)=-\frac13,\quad
  L(L)=1,\quad L(\nuc)=L(\ec)=-1.
\end{equation}
The generated P7 ledger reports
\begin{equation}
\begin{array}{c|c}
\text{quantity} & \text{value}\\
\hline
\text{allowed symbolic poles processed} & 18\\
\text{B and L conserving poles} & 12\\
\text{standard BNV seed poles} & 4\\
\text{BNV-with-sterile-neutrino poles} & 2\\
\text{baryon-violating symbolic poles} & 6\\
\text{BNV vector poles blocked by P6 completion} & 6\\
\text{physical proton bounds evaluable now} & 0\\
\text{all physical bounds conditional} & \text{true}
\end{array}
\end{equation}

The standard baryon-violating seeds include
\begin{equation}
  QQQL,\qquad \uc\uc\dc\ec,
\end{equation}
and the sterile-neutrino branch includes
\begin{equation}
  \uc\dc\dc\nuc.
\end{equation}
All of these preserve $B-L$ in the P1 bookkeeping.

The physical proton-bound interface is recorded as
\begin{equation}
  \Gamma(p\to a)
  =
  \sum_{i,j}
  C_i^{\rm phys}\,
  C_j^{\rm phys\,\ast}\,
  H_{ij}^{(a)},
\end{equation}
where $H_{ij}^{(a)}$ contains phase space, RG evolution, chiral reduction, and
hadronic matrix elements for channel $a$.  A lifetime lower bound implies
\begin{equation}
  (C^{\rm phys})^\dagger H^{(a)} C^{\rm phys}
  <
  \frac{1}{\tau_a^{\rm exp}}.
\end{equation}
For a single coefficient this reduces schematically to
\begin{equation}
  |C_i^{\rm phys}|
  <
  \sqrt{\frac{1}{\tau_a^{\rm exp}H_{ii}^{(a)}}}.
\end{equation}

P7 deliberately does not insert numerical proton limits.  Numerical
lifetimes require:
\begin{equation}
\begin{gathered}
  \text{P6-completed high-energy branch},\quad
  \text{mediator masses and couplings},\quad
  \text{physical flavor rotations},\\
  \text{operator basis and chiral contractions},\quad
  \text{RG factors},\quad
  \text{lattice/chiral matrix elements},\quad
  \tau_a^{\rm exp}.
\end{gathered}
\end{equation}
Thus P7 is a proton-bound readiness report, not a proton-lifetime
calculation.

\section{P8: Comparative Theory Scorecard}

P8 compares candidate branches using the explicit P1--P7 gates.  It is a
hard-gate ledger, not a subjective weighted score.  The gates are:
\begin{center}
\small
\begin{tabular}{
  >{\raggedright\arraybackslash}p{0.12\linewidth}
  >{\raggedright\arraybackslash}p{0.72\linewidth}}
C0 & representation pre-filter\\
P1 & charge/anomaly baseline\\
P2 & charge-filtered pole ledger\\
P3 & residue factorization/PSD\\
P4 & Ward-completion readiness\\
P5 & anomaly/chiral/exotics ledger\\
P6 & high-energy-growth completion gate\\
P7 & proton-bound readiness
\end{tabular}
\end{center}

The generated P8 scorecard compares five candidates.  Its leading conditional
finite field-theory branch is $\Spin(10)$, and it does not yet force a
tower-like UV completion.
The candidate verdicts are:
\begin{center}
\small
\begin{tabular}{
  >{\raggedright\arraybackslash}p{0.24\linewidth}
  >{\raggedright\arraybackslash}p{0.64\linewidth}}
\toprule
candidate & P8 verdict \\
\midrule
$\SU(5)$ &
conditional branch if the single-object filter is relaxed; anomaly-consistent
but family-split as $10\oplus\bar 5\oplus 1$ \\
Pati--Salam &
conditional intermediate face; natural and anomaly-consistent, but product
group and not one simple single-object branch \\
$\Spin(10)$ &
leading conditional finite field-theory branch under the current minimal
single-object filter; still has action-level and proton-bound completion debt\\
$E_6$ &
conditional branch with exotics-lifting debt from $27\to16\oplus10\oplus1$ \\
superstring-derived completions &
UV-completion class, not forced by P1--P7; serious comparison branch if
finite-pole completions fail \\
\bottomrule
\end{tabular}
\end{center}

Thus P8 closes the first Route C pass as a comparative filter:
\begin{equation}
  \Spin(10)
  =
  \text{leading conditional finite branch},
  \qquad
  \text{not a completed GUT proof}.
\end{equation}
Superstring-derived completions remain comparison branches, not P8
conclusions.

\section{Current Status and Next Check}

The P1--P8 chain currently establishes:
\begin{enumerate}
\item the exact P1 charge/anomaly baseline for one $SM+\nuc$ family;
\item the P2 charge-filtered symbolic pole components;
\item the P3 conditional residue-factorization and PSD statement;
\item the P4 Ward-identity completion ledger for broken/off-face sectors;
\item the P5 anomaly/chiral-consistency comparison across candidate branches.
\item the P6 high-energy-growth ledger for finite-vector, scalar/auxiliary,
contact-only, and tower-like comparison branches.
\item the P7 conditional low-energy matching and proton-bound readiness
ledger.
\item the P8 comparative scorecard across $\SU(5)$, Pati--Salam, $\Spin(10)$,
$E_6$, and superstring-derived completions.
\end{enumerate}
The next natural stage is no longer another broad comparison.  It is to choose
one surviving conditional branch, most naturally $\Spin(10)$, and supply the
missing action-level completion data needed by P4--P7.

\end{document}
